% Источники: 07_Программные_требования/DOC-06.3_Ранцевые_криптосистемы.md
% Генератор: docs/build_pdf.py
% Шаблон печатного представления. Содержимое подставляет docs/build_pdf.py.
\documentclass[12pt,a4paper]{article}
\usepackage{iftex}
\ifPDFTeX
  \usepackage[utf8]{inputenc}
  \usepackage[T2A]{fontenc}
  \usepackage{newtxtext,newtxmath}
\else
  \usepackage{fontspec}
  \IfFontExistsTF{Times New Roman}
    {\setmainfont{Times New Roman}}
    {\setmainfont{Libertinus Serif}}
  \IfFontExistsTF{Arial}{\setsansfont{Arial}}{\setsansfont{Latin Modern Sans}}
  \IfFontExistsTF{Menlo}{\setmonofont{Menlo}[Scale=MatchLowercase]}
    {\setmonofont{Latin Modern Mono}[Scale=MatchLowercase]}
\fi
\usepackage[russian]{babel}
\usepackage{geometry}
\geometry{left=25mm,right=15mm,top=20mm,bottom=20mm}
\usepackage{setspace}
\onehalfspacing
\usepackage{indentfirst}
\setlength{\parindent}{1.25cm}
\setlength{\parskip}{2pt}
\setlength{\emergencystretch}{3em}
\usepackage{graphicx,booktabs,longtable,array,tabularx,calc,etoolbox}
\usepackage{enumitem}
\setlist[itemize]{leftmargin=1.25cm,itemsep=2pt,topsep=4pt}
\setlist[enumerate]{leftmargin=1.25cm,itemsep=2pt,topsep=4pt}
\newcommand{\tightlist}{\setlength{\itemsep}{0pt}\setlength{\parskip}{0pt}}
\AtBeginEnvironment{longtable}{\small\singlespacing}
\setlength{\LTpre}{8pt}
\setlength{\LTpost}{8pt}
\setlength{\tabcolsep}{5pt}
\setlength{\LTleft}{0pt}
\setlength{\LTright}{0pt}
\usepackage{xurl}
\usepackage[unicode,colorlinks=true,linkcolor=black,urlcolor=blue]{hyperref}
\usepackage{bookmark}
\hypersetup{pdftitle={Ранцевые криптосистемы},pdfauthor={Георгиев Д.Д.},pdfsubject={Криптотренажёр mephi-12}}
\urlstyle{same}
\pagestyle{plain}
\setcounter{secnumdepth}{3}
\setcounter{tocdepth}{2}

\begin{document}
\hypersetup{pageanchor=false}
\begin{titlepage}
  \thispagestyle{empty}
  \centering
  {\large НИЯУ МИФИ\par}
  \vspace*{3.5cm}
  {\LARGE\bfseries Ранцевые криптосистемы\par}
  \vspace{0.5cm}
  {\large DOC-06.3\par}
  \vspace{0.6cm}
  {\large Платформа для изучения и практики курса по криптографии\par}
  {\large Криптотренажёр mephi-12\par}
  \vspace{0.7cm}
  {\large Версия 0.8\par}
  {\large Редакция от 07.11.2026\par}
  \vspace{1cm}
  \begin{tabularx}{\textwidth}{@{}p{3.4cm}Xp{4cm}@{}}
    \toprule
    Роль & ФИО & Подпись / Дата \\
    \midrule
    Разработал & Георгиев Д.Д. & \\
    \addlinespace[6pt]
    Согласовано & & \\
    \addlinespace[6pt]
    Утверждено & & \\
    \bottomrule
  \end{tabularx}
  \vfill
  {\large Москва, 2026\par}
\end{titlepage}
\hypersetup{pageanchor=true}

\section*{Лист регистрации изменений}
\addcontentsline{toc}{section}{Лист регистрации изменений}
\begingroup\small\singlespacing
\noindent
\begin{tabularx}{\textwidth}{@{}p{1.2cm}p{2cm}p{2.5cm}Xp{1.8cm}@{}}
\toprule
Версия & Дата & Автор & Содержание изменения & Утвердил \\
\midrule
0.8 & 07.11.2026 & Георгиев Д.Д. & Подготовлена первая редакция подсистемных требований; трассировка на SR-KNAP; зафиксированы известные особенности реализации. & \\
\bottomrule
\end{tabularx}
\endgroup
\clearpage
\begingroup\small\singlespacing
\tableofcontents
\endgroup
\clearpage

\hypertarget{doc-063-ранцевые-криптосистемы}{}

{\def\LTcaptype{none} % do not increment counter
\begin{longtable}[]{@{}
  >{\raggedright\arraybackslash}p{(\linewidth - 2\tabcolsep) * \real{0.3200}}
  >{\raggedright\arraybackslash}p{(\linewidth - 2\tabcolsep) * \real{0.6800}}@{}}
\toprule\noalign{}
\begin{minipage}[b]{\linewidth}\raggedright
Поле
\end{minipage} & \begin{minipage}[b]{\linewidth}\raggedright
Значение
\end{minipage} \\
\midrule\noalign{}
\endhead
\bottomrule\noalign{}
\endlastfoot
Версия & 0.8 \\
Дата & 07.11.2026 \\
Ответственный & Георгиев Д.Д. \\
Основания & Раздел \href{../04_\%D0\%A1\%D0\%B8\%D1\%81\%D1\%82\%D0\%B5\%D0\%BC\%D0\%BD\%D1\%8B\%D0\%B5_\%D1\%82\%D1\%80\%D0\%B5\%D0\%B1\%D0\%BE\%D0\%B2\%D0\%B0\%D0\%BD\%D0\%B8\%D1\%8F/DOC-05_\%D0\%A1\%D0\%B8\%D1\%81\%D1\%82\%D0\%B5\%D0\%BC\%D0\%BD\%D1\%8B\%D0\%B5_\%D1\%82\%D1\%80\%D0\%B5\%D0\%B1\%D0\%BE\%D0\%B2\%D0\%B0\%D0\%BD\%D0\%B8\%D1\%8F.pdf\#nameddest=doc-05-\%D1\%82\%D1\%80\%D0\%B5\%D0\%B1\%D0\%BE\%D0\%B2\%D0\%B0\%D0\%BD\%D0\%B8\%D1\%8F-\%D0\%BA-\%D1\%80\%D0\%B0\%D0\%BD\%D1\%86\%D0\%B5\%D0\%B2\%D1\%8B\%D0\%BC-\%D0\%B7\%D0\%B0\%D0\%B4\%D0\%B0\%D1\%87\%D0\%B0\%D0\%BC}{ранцевых задач} DOC-05; \href{DOC-06_\%D0\%9F\%D1\%80\%D0\%BE\%D0\%B3\%D1\%80\%D0\%B0\%D0\%BC\%D0\%BC\%D0\%BD\%D1\%8B\%D0\%B5_\%D1\%82\%D1\%80\%D0\%B5\%D0\%B1\%D0\%BE\%D0\%B2\%D0\%B0\%D0\%BD\%D0\%B8\%D1\%8F.pdf\#nameddest=doc-06-\%D1\%81\%D0\%B2\%D0\%BE\%D0\%B4\%D0\%BD\%D1\%8B\%D0\%B5-\%D0\%BF\%D1\%80\%D0\%BE\%D0\%B3\%D1\%80\%D0\%B0\%D0\%BC\%D0\%BC\%D0\%BD\%D1\%8B\%D0\%B5-\%D1\%82\%D1\%80\%D0\%B5\%D0\%B1\%D0\%BE\%D0\%B2\%D0\%B0\%D0\%BD\%D0\%B8\%D1\%8F}{DOC-06} \\
Плановый срок & v1.0 --- 12.11.2026 \\
\end{longtable}
}

\section{Назначение и состав}\label{1-ux43dux430ux437ux43dux430ux447ux435ux43dux438ux435-ux438-ux441ux43eux441ux442ux430ux432}

\hypertarget{1-назначение-и-состав}{}

Документ определяет программные обязанности домена ранцевых криптосистем при генерации условий, проверке решений, редакционных генерациях и хранении. Он раскрывает \href{../04_\%D0\%A1\%D0\%B8\%D1\%81\%D1\%82\%D0\%B5\%D0\%BC\%D0\%BD\%D1\%8B\%D0\%B5_\%D1\%82\%D1\%80\%D0\%B5\%D0\%B1\%D0\%BE\%D0\%B2\%D0\%B0\%D0\%BD\%D0\%B8\%D1\%8F/DOC-05_\%D0\%A1\%D0\%B8\%D1\%81\%D1\%82\%D0\%B5\%D0\%BC\%D0\%BD\%D1\%8B\%D0\%B5_\%D1\%82\%D1\%80\%D0\%B5\%D0\%B1\%D0\%BE\%D0\%B2\%D0\%B0\%D0\%BD\%D0\%B8\%D1\%8F.pdf\#nameddest=doc-05-\%D1\%82\%D1\%80\%D0\%B5\%D0\%B1\%D0\%BE\%D0\%B2\%D0\%B0\%D0\%BD\%D0\%B8\%D1\%8F-\%D0\%BA-\%D1\%80\%D0\%B0\%D0\%BD\%D1\%86\%D0\%B5\%D0\%B2\%D1\%8B\%D0\%BC-\%D0\%B7\%D0\%B0\%D0\%B4\%D0\%B0\%D1\%87\%D0\%B0\%D0\%BC}{системные требования ранцевых задач} и подчиняется общим обязательствам \href{DOC-06_\%D0\%9F\%D1\%80\%D0\%BE\%D0\%B3\%D1\%80\%D0\%B0\%D0\%BC\%D0\%BC\%D0\%BD\%D1\%8B\%D0\%B5_\%D1\%82\%D1\%80\%D0\%B5\%D0\%B1\%D0\%BE\%D0\%B2\%D0\%B0\%D0\%BD\%D0\%B8\%D1\%8F.pdf\#nameddest=doc-06-\%D1\%81\%D0\%B2\%D0\%BE\%D0\%B4\%D0\%BD\%D1\%8B\%D0\%B5-\%D0\%BF\%D1\%80\%D0\%BE\%D0\%B3\%D1\%80\%D0\%B0\%D0\%BC\%D0\%BC\%D0\%BD\%D1\%8B\%D0\%B5-\%D1\%82\%D1\%80\%D0\%B5\%D0\%B1\%D0\%BE\%D0\%B2\%D0\%B0\%D0\%BD\%D0\%B8\%D1\%8F}{DOC-06}. Математические определения ранцевой криптосистемы приведены в \href{../../02_\%D0\%A2\%D0\%B5\%D1\%85\%D0\%BD\%D0\%B8\%D1\%87\%D0\%B5\%D1\%81\%D0\%BA\%D0\%BE\%D0\%B5_\%D0\%B7\%D0\%B0\%D0\%B4\%D0\%B0\%D0\%BD\%D0\%B8\%D0\%B5/\%D0\%A2\%D0\%B5\%D1\%85\%D0\%BD\%D0\%B8\%D1\%87\%D0\%B5\%D1\%81\%D0\%BA\%D0\%BE\%D0\%B5_\%D0\%B7\%D0\%B0\%D0\%B4\%D0\%B0\%D0\%BD\%D0\%B8\%D0\%B5.pdf}{ТЗ}; здесь определяется, что делает программная часть домена.

{\def\LTcaptype{none} % do not increment counter
\begin{longtable}[]{@{}
  >{\raggedright\arraybackslash}p{(\linewidth - 2\tabcolsep) * \real{0.3200}}
  >{\raggedright\arraybackslash}p{(\linewidth - 2\tabcolsep) * \real{0.6800}}@{}}
\toprule\noalign{}
\begin{minipage}[b]{\linewidth}\raggedright
Логическая часть
\end{minipage} & \begin{minipage}[b]{\linewidth}\raggedright
Ответственность
\end{minipage} \\
\midrule\noalign{}
\endhead
\bottomrule\noalign{}
\endlastfoot
Участие в сессиях & Предоставить генерацию и проверку ранцевых задач управлению сессиями по квалификатору конфигурации \\
Генерация условия & Построить условие выбранного варианта с частично скрытыми данными и сохранить его в составе сессии \\
Проверка решения & Проверить ответ по величинам криптосистемы и вернуть вердикт с диагностикой \\
Редакционные решатели & Построить эталонную задачу по варианту или с параметрами пользователя \\
Диспетчеризация & Выбрать решатель по варианту редакционного запроса \\
Хранение & Сохранять условия сгенерированных задач между запусками \\
\end{longtable}
}

Это распределение обязанностей; оно не задаёт имена классов, таблиц или форматы внешних операций. Внешние операции сессий определены в \href{DOC-06.2_\%D0\%A1\%D0\%B5\%D1\%81\%D1\%81\%D0\%B8\%D0\%B8_\%D0\%B7\%D0\%B0\%D0\%B4\%D0\%B0\%D1\%87_\%D0\%B8_\%D0\%BF\%D0\%BB\%D0\%B0\%D0\%BD\%D0\%B8\%D1\%80\%D0\%BE\%D0\%B2\%D1\%89\%D0\%B8\%D0\%BA.pdf\#nameddest=doc-062-\%D0\%BA\%D0\%BE\%D0\%BD\%D1\%82\%D1\%80\%D0\%B0\%D0\%BA\%D1\%82\%D1\%8B-\%D1\%81\%D0\%B5\%D1\%81\%D1\%81\%D0\%B8\%D0\%B9-\%D0\%B7\%D0\%B0\%D0\%B4\%D0\%B0\%D1\%87}{контракте}; справочник эндпоинтов --- DOC-08.

\subsection{Условие и его состав}\label{ux443ux441ux43bux43eux432ux438ux435-ux438-ux435ux433ux43e-ux441ux43eux441ux442ux430ux432}

\hypertarget{условие-и-его-состав}{}

Условие ранцевой задачи --- сохраняемый JSON-объект, сопровождающий задачу сессии. Он содержит вариант, выданные фрагменты лёгкого ранца, тяжёлый ранец и данные, достаточные для восстановления остальных величин. После выдачи условие не пересчитывается из конфигурации: проверка использует его сохранённую копию, что соответствует \href{DOC-06_\%D0\%9F\%D1\%80\%D0\%BE\%D0\%B3\%D1\%80\%D0\%B0\%D0\%BC\%D0\%BC\%D0\%BD\%D1\%8B\%D0\%B5_\%D1\%82\%D1\%80\%D0\%B5\%D0\%B1\%D0\%BE\%D0\%B2\%D0\%B0\%D0\%BD\%D0\%B8\%D1\%8F.pdf\#nameddest=swr-sys-003-\%D1\%81\%D0\%BE\%D1\%81\%D1\%82\%D0\%B0\%D0\%B2-\%D1\%81\%D0\%BE\%D0\%B7\%D0\%B4\%D0\%B0\%D0\%BD\%D0\%BD\%D0\%BE\%D0\%B9-\%D1\%81\%D0\%B5\%D1\%81\%D1\%81\%D0\%B8\%D0\%B8}{составу созданной сессии}.

Приоритет всех программных требований документа --- обязательный. Источники SR указывают на определения в DOC-05; описания проверок приведены в разделе 7.

\section{Участие в сессиях}\label{2-ux443ux447ux430ux441ux442ux438ux435-ux432-ux441ux435ux441ux441ux438ux44fux445}

\hypertarget{2-участие-в-сессиях}{}

\subsection{SWR-KNAP-001. Регистрация предметного обработчика}\label{swr-knap-001-ux440ux435ux433ux438ux441ux442ux440ux430ux446ux438ux44f-ux43fux440ux435ux434ux43cux435ux442ux43dux43eux433ux43e-ux43eux431ux440ux430ux431ux43eux442ux447ux438ux43aux430}

\hypertarget{swr-knap-001-регистрация-предметного-обработчика}{}

Управление сессиями получает задачи конфигурации через предметных обработчиков. Ранцевой домен предоставляет один обработчик, известный под квалификатором «Backpack»: конфигурация называет квалификатор, а управление сессиями передаёт ему генерацию и проверку. Так состав набора задаётся без знания внутренних алгоритмов домена. Это применение \href{DOC-06_\%D0\%9F\%D1\%80\%D0\%BE\%D0\%B3\%D1\%80\%D0\%B0\%D0\%BC\%D0\%BC\%D0\%BD\%D1\%8B\%D0\%B5_\%D1\%82\%D1\%80\%D0\%B5\%D0\%B1\%D0\%BE\%D0\%B2\%D0\%B0\%D0\%BD\%D0\%B8\%D1\%8F.pdf\#nameddest=swr-sys-004-\%D0\%BA\%D0\%BE\%D0\%BD\%D1\%82\%D1\%80\%D0\%B0\%D0\%BA\%D1\%82-\%D0\%BF\%D1\%80\%D0\%B5\%D0\%B4\%D0\%BC\%D0\%B5\%D1\%82\%D0\%BD\%D0\%BE\%D0\%B3\%D0\%BE-\%D0\%BE\%D0\%B1\%D1\%80\%D0\%B0\%D0\%B1\%D0\%BE\%D1\%82\%D1\%87\%D0\%B8\%D0\%BA\%D0\%B0}{общего контракта предметного обработчика} для ранцевого типа.

\textbf{Требование.} Обработчик ранцевых задач должен быть доступен управлению сессиями под квалификатором «Backpack» и вызываться для задач конфигураций с этим квалификатором.

Вызов генерации происходит при создании набора, вызов проверки --- после допуска ответа. Обработчик не получает и не изменяет правила срока, попытки и принадлежности: они остаются у управления сессиями. Порядок вызова определён в \href{DOC-06.2_\%D0\%A1\%D0\%B5\%D1\%81\%D1\%81\%D0\%B8\%D0\%B8_\%D0\%B7\%D0\%B0\%D0\%B4\%D0\%B0\%D1\%87_\%D0\%B8_\%D0\%BF\%D0\%BB\%D0\%B0\%D0\%BD\%D0\%B8\%D1\%80\%D0\%BE\%D0\%B2\%D1\%89\%D0\%B8\%D0\%BA.pdf\#nameddest=swr-ses-004-\%D1\%84\%D0\%BE\%D1\%80\%D0\%BC\%D0\%B8\%D1\%80\%D0\%BE\%D0\%B2\%D0\%B0\%D0\%BD\%D0\%B8\%D0\%B5-\%D0\%B7\%D0\%B0\%D0\%B4\%D0\%B0\%D1\%87}{формировании задач} и \href{DOC-06.2_\%D0\%A1\%D0\%B5\%D1\%81\%D1\%81\%D0\%B8\%D0\%B8_\%D0\%B7\%D0\%B0\%D0\%B4\%D0\%B0\%D1\%87_\%D0\%B8_\%D0\%BF\%D0\%BB\%D0\%B0\%D0\%BD\%D0\%B8\%D1\%80\%D0\%BE\%D0\%B2\%D1\%89\%D0\%B8\%D0\%BA.pdf\#nameddest=swr-ses-010-\%D0\%BF\%D0\%B5\%D1\%80\%D0\%B5\%D0\%B4\%D0\%B0\%D1\%87\%D0\%B0-\%D0\%BF\%D1\%80\%D0\%BE\%D1\%84\%D0\%B8\%D0\%BB\%D1\%8C\%D0\%BD\%D0\%BE\%D0\%B9-\%D0\%BF\%D1\%80\%D0\%BE\%D0\%B2\%D0\%B5\%D1\%80\%D0\%BA\%D0\%B5}{передаче профильной проверке}.

Источники: SR-KNAP-001; SWR-SYS-004.

\subsection{SWR-KNAP-002. Случайный выбор варианта}\label{swr-knap-002-ux441ux43bux443ux447ux430ux439ux43dux44bux439-ux432ux44bux431ux43eux440-ux432ux430ux440ux438ux430ux43dux442ux430}

\hypertarget{swr-knap-002-случайный-выбор-варианта}{}

Ранцевая часть конфигурации не различает учебные варианты: оба вида заданий обслуживаются одним квалификатором. Конкретный вариант каждого задания выбирается при генерации так, чтобы состав набора не предопределял тип всех ранцевых задач. Выбор не управляется параметрами пользователя и не зависит от предыдущих задач сессии.

\textbf{Требование.} При генерации задачи для сессии ПО должно выбирать вариант случайно из двух поддерживаемых.

Выбранный вариант фиксируется в сохранённом условии и далее используется проверкой без повторного выбора. Системное основание --- \href{../04_\%D0\%A1\%D0\%B8\%D1\%81\%D1\%82\%D0\%B5\%D0\%BC\%D0\%BD\%D1\%8B\%D0\%B5_\%D1\%82\%D1\%80\%D0\%B5\%D0\%B1\%D0\%BE\%D0\%B2\%D0\%B0\%D0\%BD\%D0\%B8\%D1\%8F/DOC-05_\%D0\%A1\%D0\%B8\%D1\%81\%D1\%82\%D0\%B5\%D0\%BC\%D0\%BD\%D1\%8B\%D0\%B5_\%D1\%82\%D1\%80\%D0\%B5\%D0\%B1\%D0\%BE\%D0\%B2\%D0\%B0\%D0\%BD\%D0\%B8\%D1\%8F.pdf\#nameddest=sr-knap-001-\%D0\%B2\%D1\%8B\%D0\%B1\%D0\%BE\%D1\%80-\%D0\%B2\%D0\%B0\%D1\%80\%D0\%B8\%D0\%B0\%D0\%BD\%D1\%82\%D0\%B0-\%D1\%80\%D0\%B0\%D0\%BD\%D1\%86\%D0\%B5\%D0\%B2\%D0\%BE\%D0\%B9-\%D0\%B7\%D0\%B0\%D0\%B4\%D0\%B0\%D1\%87\%D0\%B8}{выбор варианта ранцевой задачи}.

Источник: SR-KNAP-001.

\section{Генерация условия}\label{3-ux433ux435ux43dux435ux440ux430ux446ux438ux44f-ux443ux441ux43bux43eux432ux438ux44f}

\hypertarget{3-генерация-условия}{}

\subsection{SWR-KNAP-003. Генерация сообщения}\label{swr-knap-003-ux433ux435ux43dux435ux440ux430ux446ux438ux44f-ux441ux43eux43eux431ux449ux435ux43dux438ux44f}

\hypertarget{swr-knap-003-генерация-сообщения}{}

Сообщение определяет, какие элементы тяжёлого ранца суммируются при шифровании, и входит как в условие, так и в проверяемый ответ. Учебное сообщение должно быть коротким для ручного решения и содержать достаточно нулей, чтобы расшифровка жадным алгоритмом была наглядной.

\textbf{Требование.} Сгенерированное сообщение должно быть случайной последовательностью логических значений длиной от 3 до 5 с долей ложных значений не менее трети.

Сгенерированное сообщение сохраняется в условии и сравнивается с переданным в ответе. Параметры длины относятся к задачам в сессиях; редакционные генерации принимают сообщение пользователя по SWR-KNAP-013.

Источники: SR-KNAP-001, SR-KNAP-002.

\subsection{SWR-KNAP-004. Сокрытие параметров условия}\label{swr-knap-004-ux441ux43eux43aux440ux44bux442ux438ux435-ux43fux430ux440ux430ux43cux435ux442ux440ux43eux432-ux443ux441ux43bux43eux432ux438ux44f}

\hypertarget{swr-knap-004-сокрытие-параметров-условия}{}

Выдача полного набора данных превратила бы задание в чтение ответа. Поэтому условие скрывает большую часть лёгкого ранца, оставляя открытыми отдельные элементы, и может не содержать множитель. Степень сокрытия ограничена требованием восстанавливаемости: по условию должны однозначно достраиваться все величины криптосистемы.

\textbf{Требование.} В условии открытыми должны оставаться не более двух элементов лёгкого ранца; множитель преобразования может не включаться в условие.

Сокрытие выполняется при генерации и сохраняется вместе с условием: повторное чтение сессии выдаёт те же данные. Восстановление скрытого опирается на математические свойства последовательности варианта; полная проверка достроенных величин выполняется по SWR-KNAP-009. Системное основание --- \href{../04_\%D0\%A1\%D0\%B8\%D1\%81\%D1\%82\%D0\%B5\%D0\%BC\%D0\%BD\%D1\%8B\%D0\%B5_\%D1\%82\%D1\%80\%D0\%B5\%D0\%B1\%D0\%BE\%D0\%B2\%D0\%B0\%D0\%BD\%D0\%B8\%D1\%8F/DOC-05_\%D0\%A1\%D0\%B8\%D1\%81\%D1\%82\%D0\%B5\%D0\%BC\%D0\%BD\%D1\%8B\%D0\%B5_\%D1\%82\%D1\%80\%D0\%B5\%D0\%B1\%D0\%BE\%D0\%B2\%D0\%B0\%D0\%BD\%D0\%B8\%D1\%8F.pdf\#nameddest=sr-knap-002-\%D1\%87\%D0\%B0\%D1\%81\%D1\%82\%D0\%B8\%D1\%87\%D0\%BD\%D0\%BE-\%D1\%81\%D0\%BA\%D1\%80\%D1\%8B\%D1\%82\%D0\%BE\%D0\%B5-\%D1\%83\%D1\%81\%D0\%BB\%D0\%BE\%D0\%B2\%D0\%B8\%D0\%B5}{частично скрытое условие}.

Источник: SR-KNAP-002.

\subsection{SWR-KNAP-005. Супервозрастающий лёгкий ранец}\label{swr-knap-005-ux441ux443ux43fux435ux440ux432ux43eux437ux440ux430ux441ux442ux430ux44eux449ux438ux439-ux43bux451ux433ux43aux438ux439-ux440ux430ux43dux435ux446}

\hypertarget{swr-knap-005-супервозрастающий-лёгкий-ранец}{}

Для супервозрастающего варианта лёгкий ранец обязан строго возрастать: каждый элемент больше суммы всех предыдущих. Это свойство используется и жадной расшифровкой, и проверкой ответа. Последовательность, переданная или построенная с нарушением свойства, не образует допустимую криптосистему данного варианта.

\textbf{Требование.} Для супервозрастающего варианта лёгкий ранец должен быть приведён к строго сверхвозрастающей последовательности заданной размерности.

Приведение выполняется при генерации; размерность соответствует параметрам задачи. Проверка ответа повторно устанавливает свойство для переданного ранца --- см. SWR-KNAP-009.

Источники: SR-KNAP-001, SR-KNAP-004.

\subsection{SWR-KNAP-006. Лёгкий ранец на степенях}\label{swr-knap-006-ux43bux451ux433ux43aux438ux439-ux440ux430ux43dux435ux446-ux43dux430-ux441ux442ux435ux43fux435ux43dux44fux445}

\hypertarget{swr-knap-006-лёгкий-ранец-на-степенях}{}

Вариант на степенях строит лёгкий ранец из последовательных степеней выбранного основания. Такая последовательность автоматически сверхвозрастает, что сохраняет применимость общей схемы проверки. Основание выбирается из фиксированного набора простых чисел.

\textbf{Требование.} Для варианта на степенях элементы лёгкого ранца должны строиться как степени выбранного основания.

Выбранное основание фиксируется в данных генерации и определяет ожидаемые элементы при проверке. Системное основание --- \href{../04_\%D0\%A1\%D0\%B8\%D1\%81\%D1\%82\%D0\%B5\%D0\%BC\%D0\%BD\%D1\%8B\%D0\%B5_\%D1\%82\%D1\%80\%D0\%B5\%D0\%B1\%D0\%BE\%D0\%B2\%D0\%B0\%D0\%BD\%D0\%B8\%D1\%8F/DOC-05_\%D0\%A1\%D0\%B8\%D1\%81\%D1\%82\%D0\%B5\%D0\%BC\%D0\%BD\%D1\%8B\%D0\%B5_\%D1\%82\%D1\%80\%D0\%B5\%D0\%B1\%D0\%BE\%D0\%B2\%D0\%B0\%D0\%BD\%D0\%B8\%D1\%8F.pdf\#nameddest=sr-knap-006-\%D0\%B2\%D0\%B0\%D1\%80\%D0\%B8\%D0\%B0\%D0\%BD\%D1\%82-\%D0\%BD\%D0\%B0-\%D1\%81\%D1\%82\%D0\%B5\%D0\%BF\%D0\%B5\%D0\%BD\%D1\%8F\%D1\%85}{вариант на степенях}; известное отклонение в сопоставлении элементов приведено в разделе 8.

Источники: SR-KNAP-001, SR-KNAP-006.

\subsection{SWR-KNAP-007. Сохранение сгенерированной задачи}\label{swr-knap-007-ux441ux43eux445ux440ux430ux43dux435ux43dux438ux435-ux441ux433ux435ux43dux435ux440ux438ux440ux43eux432ux430ux43dux43dux43eux439-ux437ux430ux434ux430ux447ux438}

\hypertarget{swr-knap-007-сохранение-сгенерированной-задачи}{}

Построенное условие должно стать данными создаваемой сессии, а не промежуточным значением в памяти генератора. Вместе с условием сохраняются вариант и принадлежность задачи. Отказ генерации или сохранения обрабатывается общим правилом полного набора: частичная сессия не остаётся активной.

\textbf{Требование.} Построенная задача должна сохраняться вместе с вариантом и принадлежностью, а её условие --- передаваться в состав создаваемой сессии до подтверждения операции создания.

Сохранённое условие затем используется при проверке без повторной генерации. Транзакционные границы операции создания определены \href{DOC-06.2_\%D0\%A1\%D0\%B5\%D1\%81\%D1\%81\%D0\%B8\%D0\%B8_\%D0\%B7\%D0\%B0\%D0\%B4\%D0\%B0\%D1\%87_\%D0\%B8_\%D0\%BF\%D0\%BB\%D0\%B0\%D0\%BD\%D0\%B8\%D1\%80\%D0\%BE\%D0\%B2\%D1\%89\%D0\%B8\%D0\%BA.pdf\#nameddest=swr-ses-005-\%D1\%81\%D0\%BE\%D1\%85\%D1\%80\%D0\%B0\%D0\%BD\%D0\%B5\%D0\%BD\%D0\%B8\%D0\%B5-\%D1\%81\%D0\%BE\%D0\%B7\%D0\%B4\%D0\%B0\%D0\%BD\%D0\%BD\%D0\%BE\%D0\%B3\%D0\%BE-\%D0\%BD\%D0\%B0\%D0\%B1\%D0\%BE\%D1\%80\%D0\%B0}{правилом сохранения набора} домена сессий; ранцевой обработчик участвует в ней как источник условия.

Источники: SR-KNAP-001, SR-KNAP-002; SWR-SYS-003, SWR-SES-004.

\section{Проверка решения}\label{4-ux43fux440ux43eux432ux435ux440ux43aux430-ux440ux435ux448ux435ux43dux438ux44f}

\hypertarget{4-проверка-решения}{}

\subsection{SWR-KNAP-008. Предварительная сверка ответа}\label{swr-knap-008-ux43fux440ux435ux434ux432ux430ux440ux438ux442ux435ux43bux44cux43dux430ux44f-ux441ux432ux435ux440ux43aux430-ux43eux442ux432ux435ux442ux430}

\hypertarget{swr-knap-008-предварительная-сверка-ответа}{}

До математического сравнения устанавливается, что ответ соответствует схеме ранцевого типа: передан полный набор полей, размерность --- положительное число, сообщения --- списки логических значений, целые величины --- десятичные строки, ранцы --- списки таких строк. Состав полей подтверждается настоящим документом как владелец типа в соответствии с \href{DOC-06.2_\%D0\%A1\%D0\%B5\%D1\%81\%D1\%81\%D0\%B8\%D0\%B8_\%D0\%B7\%D0\%B0\%D0\%B4\%D0\%B0\%D1\%87_\%D0\%B8_\%D0\%BF\%D0\%BB\%D0\%B0\%D0\%BD\%D0\%B8\%D1\%80\%D0\%BE\%D0\%B2\%D1\%89\%D0\%B8\%D0\%BA.pdf\#nameddest=\%D1\%80\%D0\%B0\%D0\%BD\%D1\%86\%D0\%B5\%D0\%B2\%D0\%B0\%D1\%8F-\%D0\%B7\%D0\%B0\%D0\%B4\%D0\%B0\%D1\%87\%D0\%B0}{контрактом}.

\textbf{Требование.} До предметного вердикта ПО должно сверить вариант и полноту обязательных полей ответа по схеме сохранённого варианта задачи; несоответствие должно приводить к отказу без предметной проверки.

Переданный вариант сравнивается с вариантом сохранённого условия: совпадение обязательно, подмена не допускается. Категория ошибки формы задаётся контрактом сессий (\texttt{INVALID\_REQUEST}) и не расходует попытку; далее допустимый ответ передаётся математической проверке. Разделение стадий следует \href{DOC-06.2_\%D0\%A1\%D0\%B5\%D1\%81\%D1\%81\%D0\%B8\%D0\%B8_\%D0\%B7\%D0\%B0\%D0\%B4\%D0\%B0\%D1\%87_\%D0\%B8_\%D0\%BF\%D0\%BB\%D0\%B0\%D0\%BD\%D0\%B8\%D1\%80\%D0\%BE\%D0\%B2\%D1\%89\%D0\%B8\%D0\%BA.pdf\#nameddest=swr-ses-009-\%D0\%BF\%D1\%80\%D0\%BE\%D0\%B2\%D0\%B5\%D1\%80\%D0\%BA\%D0\%B0-\%D1\%84\%D0\%BE\%D1\%80\%D0\%BC\%D0\%B0\%D1\%82\%D0\%B0}{правилу проверки формата} домена сессий.

Источники: SR-KNAP-003; SR-SES-009, SWR-SES-009, SWR-SES-022.

\subsection{SWR-KNAP-009. Пошаговая проверка величин}\label{swr-knap-009-ux43fux43eux448ux430ux433ux43eux432ux430ux44f-ux43fux440ux43eux432ux435ux440ux43aux430-ux432ux435ux43bux438ux447ux438ux43d}

\hypertarget{swr-knap-009-пошаговая-проверка-величин}{}

Предметная проверка перебирает величины криптосистемы по установленному порядку: длину и элементы сообщений и ранцев, сверхвозрастаемость, соотношение модуля с суммой лёгкого ранца, взаимную простоту множителя, тяжёлый ранец, обратный множитель, шифртекст и расшифрованное сообщение. Расшифровка выполняется жадным алгоритмом по супервозрастающей последовательности; дополнительно устанавливается положительность элементов ранцев.

\textbf{Требование.} ПО должно проверять величины, перечисленные в SR-KNAP-004, поэтапно, включая расшифровку жадным алгоритмом по супервозрастающей последовательности и положительность элементов ранцев.

Вычисление конечно и не запускает собственных фоновых процессов. Его длительность входит в \href{DOC-06.2_\%D0\%A1\%D0\%B5\%D1\%81\%D1\%81\%D0\%B8\%D0\%B8_\%D0\%B7\%D0\%B0\%D0\%B4\%D0\%B0\%D1\%87_\%D0\%B8_\%D0\%BF\%D0\%BB\%D0\%B0\%D0\%BD\%D0\%B8\%D1\%80\%D0\%BE\%D0\%B2\%D1\%89\%D0\%B8\%D0\%BA.pdf\#nameddest=swr-ses-020-\%D0\%BF\%D1\%80\%D0\%B5\%D0\%B4\%D0\%B5\%D0\%BB-\%D0\%BE\%D0\%B1\%D1\%80\%D0\%B0\%D0\%B1\%D0\%BE\%D1\%82\%D0\%BA\%D0\%B8}{предел обработки} принятого ответа: поздний результат не фиксируется, а операция завершается технической категорией без расхода попытки. Для варианта на степенях состав фактически проверяемых величин уже перечня --- см. раздел 8.

Источник: SR-KNAP-004; SR-SES-012.

\subsection{SWR-KNAP-010. Диагностика неверной величины}\label{swr-knap-010-ux434ux438ux430ux433ux43dux43eux441ux442ux438ux43aux430-ux43dux435ux432ux435ux440ux43dux43eux439-ux432ux435ux43bux438ux447ux438ux43dux44b}

\hypertarget{swr-knap-010-диагностика-неверной-величины}{}

Первая несовпавшая величина определяет текст диагностики. Пояснение относится к завершённой неверной попытке и позволяет студенту найти ошибку в собственной цепочке вычислений. Дальнейшие величины после первого несовпадения не сравниваются: вердикт уже определён.

\textbf{Требование.} При первом несовпадении ПО должно возвращать вердикт «неверно» с указанием величины, не прошедшей проверку.

Диагностика передаётся в поле пояснения предметного результата; при верном решении поле отсутствует. Текст не участвует в определении категорий отказа --- они различаются машинными кодами по \href{DOC-06.2_\%D0\%A1\%D0\%B5\%D1\%81\%D1\%81\%D0\%B8\%D0\%B8_\%D0\%B7\%D0\%B0\%D0\%B4\%D0\%B0\%D1\%87_\%D0\%B8_\%D0\%BF\%D0\%BB\%D0\%B0\%D0\%BD\%D0\%B8\%D1\%80\%D0\%BE\%D0\%B2\%D1\%89\%D0\%B8\%D0\%BA.pdf\#nameddest=\%D0\%BA\%D0\%BE\%D0\%B4\%D1\%8B-\%D0\%B8-\%D0\%BF\%D1\%80\%D0\%B5\%D0\%B4\%D1\%81\%D1\%82\%D0\%B0\%D0\%B2\%D0\%BB\%D0\%B5\%D0\%BD\%D0\%B8\%D0\%B5-\%D0\%BE\%D1\%88\%D0\%B8\%D0\%B1\%D0\%BE\%D0\%BA}{контракту ошибок}.

Источник: SR-KNAP-005.

\subsection{SWR-KNAP-011. Защита итогового состояния}\label{swr-knap-011-ux437ux430ux449ux438ux442ux430-ux438ux442ux43eux433ux43eux432ux43eux433ux43e-ux441ux43eux441ux442ux43eux44fux43dux438ux44f}

\hypertarget{swr-knap-011-защита-итогового-состояния}{}

Повторное обращение к завершённой задаче не должно запускать математическую проверку. Решение принимается по сохранённым данным задачи, поэтому состояние клиента и повторные соединения не влияют на запрет второй попытки. Восстановление отображения выполняется чтением сессии.

\textbf{Требование.} Для задачи в итоговом состоянии обработчик должен отклонять новый ответ без повторной предметной проверки, сохраняя прежний результат.

Отказ выполняется до вычисления величин. Внешняя категория определена контрактом сессий (\texttt{TASK\_FINALIZED}); порядок установлен \href{DOC-06.2_\%D0\%A1\%D0\%B5\%D1\%81\%D1\%81\%D0\%B8\%D0\%B8_\%D0\%B7\%D0\%B0\%D0\%B4\%D0\%B0\%D1\%87_\%D0\%B8_\%D0\%BF\%D0\%BB\%D0\%B0\%D0\%BD\%D0\%B8\%D1\%80\%D0\%BE\%D0\%B2\%D1\%89\%D0\%B8\%D0\%BA.pdf\#nameddest=swr-ses-013-\%D1\%81\%D0\%BE\%D1\%85\%D1\%80\%D0\%B0\%D0\%BD\%D0\%B5\%D0\%BD\%D0\%B8\%D0\%B5-\%D0\%B8\%D1\%82\%D0\%BE\%D0\%B3\%D0\%BE\%D0\%B2\%D0\%BE\%D0\%B3\%D0\%BE-\%D1\%81\%D0\%BE\%D1\%81\%D1\%82\%D0\%BE\%D1\%8F\%D0\%BD\%D0\%B8\%D1\%8F}{правилом повторного ответа} домена сессий.

Источники: SR-KNAP-007; SR-SES-008, SWR-SES-013.

\section{Редакционные генерации}\label{5-ux440ux435ux434ux430ux43aux446ux438ux43eux43dux43dux44bux435-ux433ux435ux43dux435ux440ux430ux446ux438ux438}

\hypertarget{5-редакционные-генерации}{}

\subsection{SWR-KNAP-012. Генерация по варианту}\label{swr-knap-012-ux433ux435ux43dux435ux440ux430ux446ux438ux44f-ux43fux43e-ux432ux430ux440ux438ux430ux43dux442ux443}

\hypertarget{swr-knap-012-генерация-по-варианту}{}

Редакционный запрос без дополнительных параметров выдаёт эталон выбранного варианта. Недостающие данные выбираются системой: основание --- из фиксированного набора простых чисел, остальные параметры --- по правилам варианта. Результат содержит полный набор величин и служит материалом для разбора решения.

\textbf{Требование.} Запрос выбранного варианта без дополнительных параметров должен строить эталон с полным решением; параметры, не указанные в запросе, выбираются ПО, в том числе основание --- из фиксированного набора простых чисел.

Редакционная генерация не связана с сессиями: она не проверяет допуски и не расходует попытки. Состав выдачи соответствует \href{../04_\%D0\%A1\%D0\%B8\%D1\%81\%D1\%82\%D0\%B5\%D0\%BC\%D0\%BD\%D1\%8B\%D0\%B5_\%D1\%82\%D1\%80\%D0\%B5\%D0\%B1\%D0\%BE\%D0\%B2\%D0\%B0\%D0\%BD\%D0\%B8\%D1\%8F/DOC-05_\%D0\%A1\%D0\%B8\%D1\%81\%D1\%82\%D0\%B5\%D0\%BC\%D0\%BD\%D1\%8B\%D0\%B5_\%D1\%82\%D1\%80\%D0\%B5\%D0\%B1\%D0\%BE\%D0\%B2\%D0\%B0\%D0\%BD\%D0\%B8\%D1\%8F.pdf\#nameddest=sr-knap-008-\%D1\%80\%D0\%B5\%D0\%B4\%D0\%B0\%D0\%BA\%D1\%86\%D0\%B8\%D0\%BE\%D0\%BD\%D0\%BD\%D0\%B0\%D1\%8F-\%D0\%B7\%D0\%B0\%D0\%B4\%D0\%B0\%D1\%87\%D0\%B0-\%D0\%BF\%D0\%BE-\%D0\%B2\%D0\%B0\%D1\%80\%D0\%B8\%D0\%B0\%D0\%BD\%D1\%82\%D1\%83}{эталонной задаче}.

Источник: SR-KNAP-008.

\subsection{SWR-KNAP-013. Генерация с параметрами}\label{swr-knap-013-ux433ux435ux43dux435ux440ux430ux446ux438ux44f-ux441-ux43fux430ux440ux430ux43cux435ux442ux440ux430ux43cux438}

\hypertarget{swr-knap-013-генерация-с-параметрами}{}

Параметры преподавателя делают эталон воспроизводимым в методических целях: размерность, вариант, сообщение, лёгкий ранец и множитель задаются запросом. Система проверяет математическую допустимость до построения; нарушение взаимной простоты множителя с модулем делает обратное преобразование невозможным и отклоняется.

\textbf{Требование.} Запрос с параметрами должен использовать переданные размерность, вариант, сообщение, лёгкий ранец и множитель; взаимная простота множителя с модулем должна проверяться до построения эталона.

Отклонение сопровождается диагностикой причины; частично построенный эталон не выдаётся. Отказ редакционной генерации не влияет на сессии и попытки.

Источник: SR-KNAP-009.

\subsection{SWR-KNAP-014. Диспетчеризация по варианту}\label{swr-knap-014-ux434ux438ux441ux43fux435ux442ux447ux435ux440ux438ux437ux430ux446ux438ux44f-ux43fux43e-ux432ux430ux440ux438ux430ux43dux442ux443}

\hypertarget{swr-knap-014-диспетчеризация-по-варианту}{}

Редакционные запросы двух вариантов использует общий вход, поэтому перед построением выбирается решатель соответствующего варианта. Ошибка выбора привела бы к эталону другой криптосистемы при корректно указанном варианте. Требование фиксирует обязанность; фактическое отклонение текущей реализации для варианта на степенях приведено в разделе 8, п. 1.

\textbf{Требование.} Редакционный запрос должен направляться решателю соответствующего варианта.

Выбор выполняется по варианту запроса; результат диспетчеризации --- эталон, соответствующий указанному варианту по составу величин.

Источники: SR-KNAP-008, SR-KNAP-009.

\subsection{SWR-KNAP-015. Генерация без сохранения}\label{swr-knap-015-ux433ux435ux43dux435ux440ux430ux446ux438ux44f-ux431ux435ux437-ux441ux43eux445ux440ux430ux43dux435ux43dux438ux44f}

\hypertarget{swr-knap-015-генерация-без-сохранения}{}

Редакционные примеры не являются попытками и не должны попадать в хранилище сессий. Результат строится на лету и возвращается напрямую; повторный запрос выполняет генерацию заново.

\textbf{Требование.} Результаты редакционных генераций не должны записываться в хранилище.

Правило не затрагивает сохранение задач, создаваемых в сессиях по SWR-KNAP-007.

Источник: SR-KNAP-010.

\section{Хранение}\label{6-ux445ux440ux430ux43dux435ux43dux438ux435}

\hypertarget{6-хранение}{}

\subsection{SWR-KNAP-016. Сохранность условий}\label{swr-knap-016-ux441ux43eux445ux440ux430ux43dux43dux43eux441ux442ux44c-ux443ux441ux43bux43eux432ux438ux439}

\hypertarget{swr-knap-016-сохранность-условий}{}

Условия сгенерированных ранцевых задач должны переживать перезапуск приложения: вместе с вариантом и принадлежностью они образуют данные, по которым выполняется проверка и чтение результата. Списковые величины хранятся в текстовом представлении колонок и восстанавливаются без изменения значений.

\textbf{Требование.} Условия сгенерированных ранцевых задач должны сохраняться между запусками приложения вместе с вариантом и принадлежностью; списковые величины --- в текстовом представлении; сохранённые условия доступны для просмотра истории.

Сохранённые условия читаются историей задач без повторной генерации. Общие границы сохранения и восстановления определены \href{DOC-06_\%D0\%9F\%D1\%80\%D0\%BE\%D0\%B3\%D1\%80\%D0\%B0\%D0\%BC\%D0\%BC\%D0\%BD\%D1\%8B\%D0\%B5_\%D1\%82\%D1\%80\%D0\%B5\%D0\%B1\%D0\%BE\%D0\%B2\%D0\%B0\%D0\%BD\%D0\%B8\%D1\%8F.pdf\#nameddest=swr-sys-008-\%D1\%81\%D0\%BE\%D1\%85\%D1\%80\%D0\%B0\%D0\%BD\%D0\%BD\%D0\%BE\%D1\%81\%D1\%82\%D1\%8C-\%D0\%B7\%D0\%B0\%D0\%B2\%D0\%B5\%D1\%80\%D1\%88\%D1\%91\%D0\%BD\%D0\%BD\%D1\%8B\%D1\%85-\%D0\%B4\%D0\%B0\%D0\%BD\%D0\%BD\%D1\%8B\%D1\%85}{правилом сохранности} DOC-06.

Источники: SR-SYS-007; SWR-SYS-008.

\section{Трассировка и описательные проверки}\label{7-ux442ux440ux430ux441ux441ux438ux440ux43eux432ux43aux430-ux438-ux43eux43fux438ux441ux430ux442ux435ux43bux44cux43dux44bux435-ux43fux440ux43eux432ux435ux440ux43aux438}

\hypertarget{7-трассировка-и-описательные-проверки}{}

Ниже перечислены связи всех программных требований документа. Это описания проверки требований, не результаты прогонов и не выпущенные кейсы DOC-10.

{\def\LTcaptype{none} % do not increment counter
\begin{longtable}[]{@{}
  >{\raggedright\arraybackslash}p{(\linewidth - 4\tabcolsep) * \real{0.2512}}
  >{\raggedright\arraybackslash}p{(\linewidth - 4\tabcolsep) * \real{0.2512}}
  >{\raggedright\arraybackslash}p{(\linewidth - 4\tabcolsep) * \real{0.4976}}@{}}
\toprule\noalign{}
\begin{minipage}[b]{\linewidth}\raggedright
SWR
\end{minipage} & \begin{minipage}[b]{\linewidth}\raggedright
Системные источники
\end{minipage} & \begin{minipage}[b]{\linewidth}\raggedright
Условия и ожидаемый результат проверки
\end{minipage} \\
\midrule\noalign{}
\endhead
\bottomrule\noalign{}
\endlastfoot
SWR-KNAP-001 & SR-KNAP-001 & Задачи конфигурации с ранцевым квалификатором создаются этим обработчиком; другие типы ему не передаются \\
SWR-KNAP-002 & SR-KNAP-001 & В сессиях появляются оба варианта; пользователь не влияет на выбор \\
SWR-KNAP-003 & SR-KNAP-001, SR-KNAP-002 & Сообщения выданных условий имеют длину 3--5 и достаточную долю ложных значений \\
SWR-KNAP-004 & SR-KNAP-002 & В условии скрыты все элементы ранца, кроме не более двух; множитель может отсутствовать \\
SWR-KNAP-005 & SR-KNAP-001, SR-KNAP-004 & Лёгкий ранец супервозрастающего варианта выдерживает проверку сверхвозрастаемости \\
SWR-KNAP-006 & SR-KNAP-001, SR-KNAP-006 & Элементы ранца варианта на степенях соответствуют степеням основания \\
SWR-KNAP-007 & SR-KNAP-001, SR-KNAP-002 & Созданная сессия содержит сохранённые условия с вариантами и принадлежностью \\
SWR-KNAP-008 & SR-KNAP-003, SR-SES-009 & Ответ с чужим вариантом или без обязательных полей отклоняется до вердикта \\
SWR-KNAP-009 & SR-KNAP-004, SR-SES-012 & Искажение любой величины даёт неверный вердикт; вычисление укладывается в предел обработки \\
SWR-KNAP-010 & SR-KNAP-005 & Диагностика указывает первую несовпавшую величину \\
SWR-KNAP-011 & SR-KNAP-007, SR-SES-008 & Повторный ответ после итогового состояния не создаёт второго вердикта \\
SWR-KNAP-012 & SR-KNAP-008 & Эталон содержит полный набор величин выбранного варианта \\
SWR-KNAP-013 & SR-KNAP-009 & Переданные параметры используются; невзаимно простой множитель --- отказ с причиной \\
SWR-KNAP-014 & SR-KNAP-008, SR-KNAP-009 & Запрос направляется решателю варианта; отклонение отмечено в разделе 8 \\
SWR-KNAP-015 & SR-KNAP-010 & После редакционной генерации хранилище не пополняется \\
SWR-KNAP-016 & SR-SYS-007 & После перезапуска условия, варианты и принадлежность доступны \\
\end{longtable}
}

\section{Известные особенности}\label{8-ux438ux437ux432ux435ux441ux442ux43dux44bux435-ux43eux441ux43eux431ux435ux43dux43dux43eux441ux442ux438}

\hypertarget{8-известные-особенности}{}

Фактическое поведение реализации, расходящееся с требованиями выше или требующее командного решения. Каждая позиция --- кандидат DEF на формальной инспекции DOC-06.3 или вопрос к версии 1.0.

\begin{enumerate}
\def\labelenumi{\arabic{enumi}.}
\tightlist
\item
  В диспетчеризации редакционных запросов для варианта на степенях фактически вызывается решатель супервозрастающего варианта; ветвь решателя степеней недостижима. SWR-KNAP-014 фиксирует обязанность, расхождение --- предмет решения к v1.0.
\item
  Для варианта на степенях расшифрованное сообщение не проверяется (в коде отмечено как незавершённое). SR-KNAP-006 описывает схему целиком; фактический состав проверок уже --- предмет уточнения к v1.0.
\item
  Приведение лёгкого ранца варианта на степенях сравнивает элементы со строковым представлением степеней основания.
\item
  Поле варианта в предметном ответе передаётся как наименование на русском языке; сверка варианта идёт по наименованию, а не по коду перечисления. Отражено в \href{DOC-06.2_\%D0\%A1\%D0\%B5\%D1\%81\%D1\%81\%D0\%B8\%D0\%B8_\%D0\%B7\%D0\%B0\%D0\%B4\%D0\%B0\%D1\%87_\%D0\%B8_\%D0\%BF\%D0\%BB\%D0\%B0\%D0\%BD\%D0\%B8\%D1\%80\%D0\%BE\%D0\%B2\%D1\%89\%D0\%B8\%D0\%BA.pdf\#nameddest=\%D1\%80\%D0\%B0\%D0\%BD\%D1\%86\%D0\%B5\%D0\%B2\%D0\%B0\%D1\%8F-\%D0\%B7\%D0\%B0\%D0\%B4\%D0\%B0\%D1\%87\%D0\%B0}{контракте сессий} как строка варианта и подлежит отражению в DOC-08.
\item
  Редакционный контроллер помечен в коде как устаревший, но остаётся доступным; судьба его операций --- вопрос к v1.0.
\item
  Арифметика (наибольший общий делитель, обратный элемент) реализована внутри домена; математическая библиотека домена Эль-Гамаля не используется. При чтении документов домены не смешиваются.
\end{enumerate}

\section{Условия следующей версии}\label{9-ux443ux441ux43bux43eux432ux438ux44f-ux441ux43bux435ux434ux443ux44eux449ux435ux439-ux432ux435ux440ux441ux438ux438}

\hypertarget{9-условия-следующей-версии}{}

До версии 1.0 требуется подтвердить обязательность полей предметного ответа по вариантам в \href{DOC-06.2_\%D0\%A1\%D0\%B5\%D1\%81\%D1\%81\%D0\%B8\%D0\%B8_\%D0\%B7\%D0\%B0\%D0\%B4\%D0\%B0\%D1\%87_\%D0\%B8_\%D0\%BF\%D0\%BB\%D0\%B0\%D0\%BD\%D0\%B8\%D1\%80\%D0\%BE\%D0\%B2\%D1\%89\%D0\%B8\%D0\%BA.pdf\#nameddest=\%D1\%80\%D0\%B0\%D0\%BD\%D1\%86\%D0\%B5\%D0\%B2\%D0\%B0\%D1\%8F-\%D0\%B7\%D0\%B0\%D0\%B4\%D0\%B0\%D1\%87\%D0\%B0}{контракте сессий} (пометку «для согласования с Георгиевым» снимает редактор), согласовать передачу больших целых чисел с владельцами DOC-08 и решить судьбу устаревшего редакционного контроллера и отклонения из п. 8.1. Численные пределы размерности и длительности генерации не устанавливаются до решения владельца конфигураций.

История: 07.11.2026 --- подготовлена версия 0.8: обязанности домена, трассировка на SR-KNAP, известные особенности реализации.

Редакция 07.11.2026: развёрнуты пояснения каждого требования, добавлены примеры и точные ссылки на связанные разделы. Формулировки обязательств и идентификаторы сохранены.


\end{document}
